\hypertarget{3-methodology}{%
\section{Methodology}\label{3-methodology}}

\hypertarget{31-mfsc-discrete-event-simulation}{%
\subsection{MFSC discrete-event
simulation}\label{31-mfsc-discrete-event-simulation}}

The experimental platform is a 13-operation military food supply chain
reconstructed in Python/SimPy from Garrido-Rios' thesis~\cite{garrido2017}
and checked against the original Garrido workbook artifacts used for
raw order ledgers and APj/RPj/DPj/Re summaries. The model preserves the
original operational structure rather than replacing it with a stylized
benchmark environment. This matters for the paper's positioning: the
contribution is not a generic RL sandbox, but a thesis-grounded
operational model whose buffers, transport stages, and disruption
processes already carry domain meaning.

The DES retains hourly granularity where the original supply chain
requires sub-day behavior, especially around the assembly-line
operations. This avoids collapsing operational effects that would
disappear under daily aggregation, including short breakdowns and
shift-sensitive capacity changes. The warm-up regime also follows the
validated thesis lineage, with the first full batch reaching the
downstream chain after roughly \texttt{838.8} simulated hours. The RL
layer therefore sits on top of a model that is already operationally
constrained before learning begins.

The simulation backbone is not treated as a moving target. Deterministic
and stochastic checks remain anchored to the thesis reference tables,
and the benchmark story is built on the audited post-alignment DES
rather than on earlier historical artifacts. This allows the paper to
claim benchmark fidelity without claiming that the model is a universal
representation of all supply chains.

\begin{figure}[htbp]
\centering
\includegraphics[width=\linewidth]{fig2_mfsc_topology}
\caption{Simplified MFSC topology redrawn from Garrido-Rios' 13-operation military food supply chain. Track B preserves the DES backbone but exposes downstream dispatch controls at Op10 and Op12, the operational recovery valves left fixed in Track A.}
\label{fig:mfsc_topology}
\end{figure}

\hypertarget{32-control-formulation}{%
\subsection{Control formulation}\label{32-control-formulation}}

The learning problem is defined as weekly operational control over the
DES. At each decision epoch, the agent observes a state vector derived
from the simulation, selects a continuous action, and receives a
resilience-oriented reward. Decisions are taken every \texttt{168}
simulated hours, which matches the weekly planning cadence already
embedded in procurement, disruption, and operating rhythms. Formally, we
treat the simulator-control loop as a finite-horizon partially observed
control problem
\[
\mathcal{M}=(\mathcal{S},\mathcal{O},\mathcal{A},P,r,\gamma,H),
\]
where $s_t\in\mathcal{S}$ is the full DES state (queues, buffers,
orders, disruptions, and resource states), $o_t\in\mathcal{O}$ is the
feature vector exposed to the policy, $a_t\in\mathcal{A}$ is the
weekly operational-control action, $P$ is the simulator transition
kernel induced by SimPy events and stochastic disruptions, $r$ is the
training reward, $\gamma$ is the discount factor, and $H$ is the
evaluation horizon (52 or 104 weeks in the generalization study). The
paper-facing outcome is not raw training reward but a trajectory-level
metric panel $G(\tau)$ headed by Garrido/Excel ReT.

The paper uses two control contracts. \texttt{Track\ A} is the
thesis-grounded benchmark family. It uses observation version
\texttt{v1} with \texttt{15} dimensions and a \texttt{5D} action
contract covering upstream quantity and reorder multipliers together
with assembly-shift selection. This is the negative paper lane: it asks
whether RL can beat strong static policies when control remains limited
to the original scope.

\texttt{Track\ B} is a separate research lane rather than a retrofit of
\texttt{Track\ A}. It uses observation version \texttt{v7} with
\texttt{52} dimensions and the \texttt{track\_b\_v1} action contract
(\texttt{8D}). The action contract keeps the Track A controls and adds
downstream dispatch multipliers at \texttt{Op10} and \texttt{Op12}. The
paper's central question is whether this operational extension changes
the benchmark outcome by exposing the operational bottleneck that Track
A leaves mostly outside the controllable interface.

Operationally, the Op10/Op12 dispatch multipliers represent short-term
downstream transport expediting authority --- additional dispatch
frequency, vehicle allocation, or priority release at the two
line-of-communication stages feeding the combat service support units and
the theatre of operations. They do not alter demand, risk realization, or
upstream production physics; Track B preserves the DES process logic and
risk structure while extending the controllable decision interface.
Because the Garrido-Rios source does not provide calibrated transport
cost data, dispatch costs are treated through an explicit sensitivity
analysis (Section~4.3) rather than as calibrated monetary costs.

\begin{table}[htbp]
\centering
\caption{Observation classes relevant to the privileged-observation audit.}
\label{tab:observation_classes}
\begin{tabular}{p{0.25\linewidth}p{0.32\linewidth}p{0.30\linewidth}}
\toprule
Class & Examples & Manuscript role \\
\midrule
Operational feedback & backlog, queues, rolling fill, inventory pressure & Plausibly measurable state for closed-loop recovery \\
Derived simulator summaries & smoothed service and risk-pressure features & Reported as part of the benchmark contract \\
Privileged fields & five true-regime one-hot fields plus two transition-matrix forecast fields & Removed in E2; also tested by the E1 regime-table comparator \\
\bottomrule
\end{tabular}
\end{table}

\hypertarget{33-reward-and-reporting-metrics}{%
\subsection{Reward and reporting
metrics}\label{33-reward-and-reporting-metrics}}

The canonical Track B training reward is \texttt{control\_v1}, a linear
operational penalty (backorder fraction plus a small shift-switching
cost) rather than a resilience-index reward. This is a deliberate
design choice, not an oversight: the agent is trained on an operational
proxy and evaluated on the untouched Garrido/Excel resilience metric
(\texttt{ReT}), so the training signal cannot be tuned to directly game
the reported evaluation metric. Track A results draw on a broader set of
resilience-aligned reward families (including Cobb--Douglas and
sequential formulations) tested across the dense-frontier evaluation; no
tested reward family closes the Track A gap.

The paper-facing comparison prioritizes service and resilience outcomes
rather than raw reward totals across incompatible reward families. The
primary metric for the headline claim is \texttt{order\_ret\_excel}
(Garrido/Excel ReT), with \texttt{order\_level\_ret\_mean},
\texttt{flow\_fill\_rate}, \texttt{terminal\_rolling\_fill\_rate\_4w},
recovery-tail metrics (\texttt{order\_ctj\_p99},
\texttt{order\_rpj\_p99}), service-loss AUC, and a shift-utilization cost
index (an assembly-labor proxy; it does not price dispatch, which is
costed separately through the sensitivity in Section~4.3)
reported as a secondary evidence panel, not substitutes for the primary
metric. \texttt{reward\_total} is only comparable within a shared reward
family; comparisons across Track A and Track B, or across reward
families within Track A, rely on the evaluation metric panel rather than
on raw cumulative reward.

\begin{figure}[htbp]
\centering
\includegraphics[width=\linewidth]{fig7_ret_metric_lineage}
\caption{Garrido/Excel ReT evaluation lineage redrawn from Garrido-Rios' resilience metric structure. The paper reports order-level ReT as the primary outcome, alongside service and recovery-tail metrics, while PPO is trained on a separate operational proxy reward.}
\label{fig:ret_metric_lineage}
\end{figure}

\hypertarget{34-experimental-protocol}{%
\subsection{Experimental protocol}\label{34-experimental-protocol}}

The experimental design contains two benchmark families and one matched
ablation. \texttt{Track\ A} is the thesis-grounded benchmark lane. It
retains the validated DES backbone, the original operational scope, and
strong static baselines as the main comparators. Learned policies are
reported honestly as negative evidence when they remain below the
strongest static baseline; this includes PPO in the frozen
\texttt{ReT\_seq\_v1} lane and RecurrentPPO in the corresponding
recurrent comparator lane.

\texttt{Track\ B} is the positive paper lane. Its canonical contract fixes
the reward (\texttt{control\_v1}), the observation version (\texttt{v7},
52 dimensions), the action contract (\texttt{track\_b\_v1}, 8 dimensions),
the risk level (\texttt{adaptive\_benchmark\_v2}), the thesis year basis,
a 168-hour decision step, and stochastic processing times. The confirmatory benchmark is trained for
\texttt{60{,}000} timesteps per seed with seeds \texttt{1--5} and
evaluated over \texttt{12} common-random-number episodes per seed against
a dense downstream-dispatch static grid of 147 comparators (3 shift
levels $\times$
7 Op10 multipliers $\times$ 7 Op12 multipliers, each in $\{0.5, 0.75,
1.0, 1.25, 1.5, 2.0, 2.5\}$). The grid is dense over the dispatch/shift
family that Track B is designed to test; the remaining upstream
buffer/reorder dimensions are held at their canonical Track B settings
for all static comparators, and learned control over that upstream family
alone is separately bounded by the Track A dense-frontier negative result
(Section~4.2). As a local static-comparator bound, we also evaluate a
3$\times$3 grid over Op3 and Op9 quantity multipliers
(\{0.75, 1.00, 1.25\}) at the best downstream-dispatch cell
(S2, Op10$\times$2.00, Op12$\times$1.50) under the same CRN protocol.
This bound is not an exhaustive eight-dimensional frontier, but it tests
whether the nearest upstream static degrees of freedom overturn the
downstream-grid result.
The learner is PPO with a two-hidden-layer
(64$\times$64) multilayer-perceptron policy and value network
(\texttt{MlpPolicy}, separate actor/critic heads), learning rate
$3\times10^{-4}$, $n_{\mathrm{steps}}{=}1024$, batch size 256, 10 epochs
per update, $\gamma{=}0.99$, GAE $\lambda{=}0.95$, and clip range 0.2.

\subsection{Statistical inference}\label{35-statistical-inference}

All learned-versus-static comparisons are paired under common random
numbers~\cite{law2015simulation}: every policy is evaluated on identical
disruption and demand realizations, so policy differences are isolated
from scenario noise. The promoted Track B headline uses 120 evaluation
pairs (10 seeds $\times$ 12 episodes); the E1-E4 mechanism checks and
top-12 appendix use their original 5-seed paired designs. Episode-level
pooled confidence intervals are reported as descriptive only. Primary
inference treats the \emph{training seed} as the inferential unit:
per-seed mean paired deltas are computed first, and claims rest on seed
sign consistency together with seed-clustered bootstrap or $t$-based
intervals. To guard against winner's-curse comparator
selection, the headline delta is additionally recomputed against each of
the twelve strongest static comparators individually; all twelve deltas
remain positive with confidence intervals excluding zero
(Appendix~\ref{appendix-robustness}). Secondary metrics (fill, tails,
cost) are reported as an evidence panel without multiplicity adjustment
and are not used as stand-alone claims.

To separate the action-contract argument from action-space size, the
paper uses a matched action-space ablation at the same 5-seed,
60{,}000-timestep scale, comparing the full \texttt{joint} 8D contract
against \texttt{downstream\_only} (shift frozen) and \texttt{shift\_only}
(dispatch frozen) variants, with observation version, reward family,
risk profile, and training budget held fixed. This decomposition provides
mechanism evidence that downstream control access, not action-space
size, drives the Track B gain; it is supportive evidence inside the
studied benchmark family, not causal proof outside
it.

\begin{table}[htbp]
\centering
\caption{Experimental evidence map used in the manuscript.}
\label{tab:evidence_map}
\begin{tabular}{p{0.18\linewidth}p{0.38\linewidth}p{0.30\linewidth}}
\toprule
Lane & Purpose & Claim supported \\
\midrule
Track A & Original buffer/shift control against dense static frontiers & Boundary/null: original action surface does not yield a publishable dynamic-control win \\
Track B canonical & PPO vs 147 dense static dispatch policies under CRN & Main adaptive-recovery result \\
E1 & Static, heuristic, and true-regime lookup-table comparators & PPO win is not reproduced by zero-learning privileged regime access \\
E2 & PPO retrained with privileged regime/forecast fields masked & PPO win does not require oracle observation fields \\
E3 & Frozen PPO across risk levels and horizons & Cross-regime stress transfer; severe service-floor boundary \\
E4 & Joint, downstream-only, and shift-only ablation & Mechanism evidence for downstream bottleneck authority \\
\bottomrule
\end{tabular}
\end{table}

All DES code, Gymnasium wrappers, training and evaluation scripts, static
frontier definitions, and the frozen evidence bundles behind every table
in this paper are available in the project repository
(\texttt{github.com/Thom-320/scres-ia}, branch
\texttt{codex/garrido-replication-experiments}); each reported number
traces to a dated artifact directory and the training seeds are stated
explicitly throughout.
